\section{Methods to correct the bias of Tanimoto dot product random features}
\label{apdx:bias-correction}

Here, we describe the bias correction strategies mentioned in Section \ref{sec:tdp-rf} and tested experimentally in section~\ref{sec:experiments}.
We begin by noting that when $x,x'\in\R_{\ge 0}^d$, $x\cdot x'\ge 0$, and therefore, the power series
\eq{
T_{DP}(x,x')=\sum_{r=1}^\infty (x\cdot x')^r(\|x\|^2+\|x'\|^2)^{-r}
}
is monotone. Therefore, if we use an unbiased sketch for the truncated series 
\eq{
\sum_{r=1}^R (x\cdot x')^r(\|x\|^2+\|x'\|^2)^{-r}
}
as the one constructed in Section \ref{sec:tdp-rf}, the final sketch will be biased \emph{downward}:
\eq{
\E(\Phi(x)\cdot\Phi(x'))< T_{DP}(x,x')
}
for all $x,x'\in\R_{\ge 0}^d$. Below we introduce two strategies to remedy this issue.

\subsection{Bias correction strategy 1: normalize the features}

Empirically, we observe that the highest bias occurs in the diagonal elements of the kernel matrix. As the kernel satisfies $T_{DP}(x,x)=1$ for all $x\in\R^d$, one possible correction is to \emph{normalize} the sketch to obtain
\eq{
\tilde \Phi(x) = \frac{\Phi(x)}{\|\Phi(x)\|}.
}
This remains an oblivious sketch, and ensures that the diagonal of the kernel matrix is estimated exactly, at the expense of possibly introducing some bias in off-diagonal elements of $K$.

\subsection{Bias correction strategy 2: sketch the residual}

To simplify the algebra,
let $t_{x,y}=\frac{x\cdot y}{\|x\|^2 + \|y\|^2}$.
The power series for $T_{DP}$ can then be re-written as:
\begin{align}
    T_{DP}(x,y) &= \sum_{r=1}^\infty \left(t_{x,y}\right)^r \nonumber\\
    &= \sum_{r=1}^{R} \left(t_{x,y}\right)^r + \left(t_{x,y}\right)^{R+1} + \sum_{r=R+2}^{\infty} \left(t_{x,y}\right)^r \nonumber\\
    &= \sum_{r=1}^{R} \left(t_{x,y}\right)^r + \left(t_{x,y}\right)^{R+1} + \left(t_{x,y}\right)^{R+1}\sum_{r=1}^{\infty} \left(t_{x,y}\right)^r \nonumber\\
    &= \underbrace{\sum_{r=1}^{R} \left(t_{x,y}\right)^r}_{k^{R}} + \underbrace{\left(t_{x,y}\right)^{R+1} \left(1 + T_{DP}(x,y)\right)}_{\text{truncation error}}
\end{align}
Critically, the truncation error can be written in terms of the kernel value itself, without any infinite sums.
This enables a simple procedure to generate random features for \emph{both} the truncated
power series and the remainder:
\begin{enumerate}
    \item Compute and store $\Phi(x)$ 
        as random features for the truncated kernel $k^{R}$.
    \item  Concatenate a single 1 onto $\Phi(x)$
        to produce features $\Phi_{+1}(x) = (1,\Phi(x))$ which approximate the kernel $1+k^{R}$
        using only a single extra dimension.
        Treat this as approximate random features for the kernel $1+k$.
    \item Compute random features $\Phi_{r+1}(x)$ for $t_{x,y}^{R+1}$.
    \item Apply \textsc{Sketch} to the tensor product $\Phi_{r+1}(x)\otimes \Phi_{+1}(x)$ to obtain $\Delta(x)$,
        which approximates random features of the truncation error 
        $t_{x,y}^{R+1}(1+T_{DP}(x,y))$.
    \item Concatenate the random features $\Phi(x)$ and $\Delta(x)$, to obtain bias corrected random features $\Phi_\text{bc}(x)=\Phi\oplus \Delta(x)$.
\end{enumerate}

Overall, these random features are essentially a concatenation of
the random features for the truncated power series
with an additional random feature estimate of the truncation error,
which is formed using both the random features for the first $R$ terms
and the random features for the $(R+1)$th term.
A nice property of the procedure above is that it \emph{re-uses}
the random features $\Phi(x)$ to estimate the error.